\section{Poisson 过程与独立增量}
\label{sec:06-Random-Processes/02-Counting-Processes/02}

离散时间的 Bernoulli 模型有个硬伤：它要求事件恰好落在你画好的格点上。但电话可以在 3.72 秒时响起，粒子可以在任意时刻衰变，公交车不会只在你打卡的瞬间到站。你需要一个连续时间模型。

\subsection{从离散到连续：把网格无限加密}

回到 Bernoulli 过程的思路：把时间切成等长小段，每段独立，概率与段长成正比。现在把长度 \(t\) 的区间切成 \(n\) 段，每段长 \(t/n\)。若事件概率与段长成比例，每段发生一次事件的概率约为

\[
\lambda\frac{t}{n},
\]

多次事件的概率应该比单次事件小得多——在足够短的时间内，发生两次以上的可能性可以忽略。令 \(n\to\infty\)，段长趋于零，离散逼近变为连续极限。

固定 \(k\)，前 \(n\) 段中恰有 \(k\) 段发生事件的概率为

\[
\binom{n}{k}\left(\lambda\frac{t}{n}\right)^k
\left(1-\lambda\frac{t}{n}\right)^{\,n-k}.
\]

取 \(n\to\infty\)：

\[
\binom{n}{k}\left(\lambda\frac{t}{n}\right)^k
=\frac{n(n-1)\cdots(n-k+1)}{k!}\cdot\frac{(\lambda t)^k}{n^k}
\to\frac{(\lambda t)^k}{k!},
\]
\[
\left(1-\lambda\frac{t}{n}\right)^{n-k}
\to e^{-\lambda t}.
\]

两式相乘得到

\[
P(N(t)=k)
=e^{-\lambda t}\frac{(\lambda t)^k}{k!}.
\]

这就是 Poisson 分布——二项分布在事件概率趋于零、试验次数趋于无穷时的极限。它是处理``大量独立罕见事件''的自然分布：每个小时隙碰撞一次的概率极低，但时隙极多，结果是有限且可变的计数。

\subsection{四条公理：用条件替代极限推导}

建立 Poisson 过程的正式定义不依赖极限论证。一个速率为 \(\lambda>0\) 的齐次 Poisson 过程 \(\{N(t):t\ge0\}\) 满足四条公理：

\begin{enumerate}
\item
  \(N(0)=0\)——从零开始计数；
\item
  不重叠区间的增量独立——知道过去的事件数量不能告诉你未来会发生多少；
\item
  长度相同的区间具有相同增量分布——统计行为不随时间平移而改变；
\item
  长度 \(h\) 的极短区间中

  \[
  P(N(h)=1)=\lambda h+o(h),
  \]

  \[
  P(N(h)\ge2)=o(h).
  \]
\end{enumerate}

最后一条需要解释。\(o(h)\) 表示除以 \(h\) 后极限为零的项——它不仅小，而且比 \(h\) 小得更快。\(P(N(h)\ge2)=o(h)\) 排除了短时间内发生多次事件的可能：两次以上事件的概率不仅随 \(h\to0\) 趋于零，而且比 \(h\) 本身趋于零更快。这意味着在极小尺度上，过程至多跳一次。

这四条公理完全决定过程的分布。更一般地，对任意 \(0\le s<t\)：

\[
N(t)-N(s)
\sim\mathrm{Poisson}(\lambda(t-s)).
\]

因此

\[
\E[N(t)]
=\Var(N(t))
=\lambda t.
\]

均值和方差相等是 Poisson 分布的标志性特征。如果一个计数数据的方差显著大于均值（过度离散），Poisson 模型就需要重新审视。

\subsection{路径与小区间假设}

Poisson 过程的样本路径是右连续的整数阶梯函数：从 0 出发，在事件发生时刻向上跳 1，在跳跃之间保持常数。条件 \(P(N(h)\ge2)=o(h)\) 保证了每次跳跃幅度恰好为 1——没有同一瞬间发生两次事件的``成批跳跃''。

还有一个容易被忽略的条件：任意有限区间内事件数有限。否则过程可能在有限时间内爆炸——累积无穷多次跳跃。标准 Poisson 过程不会爆炸，因为任意有限 \(t\) 的计数 \(N(t)\) 是有限 Poisson 随机变量，以概率 1 取有限值。但在更一般点过程中，爆炸是真实的可能性，需要显式排除。

\subsection{叠加：把两条独立流合并}

设想两个独立的 Poisson 过程 \(N_1,N_2\)，速率分别为 \(\lambda_1,\lambda_2\)。它们的叠加

\[
N(t)=N_1(t)+N_2(t)
\]

还是 Poisson 过程吗？速率应该是 \(\lambda_1+\lambda_2\)——两个独立事件流的总频率就是各自频率之和。

这个直觉是对的，但需要验证的不只是单个时刻的分布。固定 \(t\)，两个独立 Poisson 之和当然是 Poisson\((\lambda_1+\lambda_2)\)。但过程级别的结论更强：叠加后的增量仍然独立、仍然平稳。需要检查任意不相交区间上的 \(N_1\) 增量与 \(N_2\) 增量组合后仍保持独立性——这来自两个过程各自的独立增量以及它们彼此独立。过程级结论比单点分布更强。

\subsection{稀疏化：随机筛选后还是 Poisson}

反过来，从一个 Poisson 流中随机筛选。对速率 \(\lambda\) 的 Poisson 过程，每个到达独立以概率 \(p\) 标记为 A 类，否则为 B 类。那么 A 类流和 B 类流是速率分别为

\[
p\lambda,\qquad(1-p)\lambda
\]

的独立 Poisson 过程。

验证思路：固定 \(t\)，条件于总数 \(N(t)=n\) 时，A 类事件数服从 \(\mathrm{Binomial}(n,p)\)。对总数边缘化：

\[
P(N_A(t)=k,N_B(t)=m)
=P(N(t)=k+m)\binom{k+m}{k}p^k(1-p)^m.
\]

代入 Poisson 质量函数，二项系数与 Poisson 质量函数的乘积恰好分解为两个 Poisson 质量函数之积——速率分别为 \(p\lambda\) 和 \((1-p)\lambda\)。这不仅给出了两个边缘分布，同时证明了固定 \(t\) 处两个计数的独立性。再结合增量独立，两个筛选后的流作为随机过程独立。

这个性质在应用中非常有用：电信路由（将呼叫随机分配到不同线路）、粒子检测（探测器对每个粒子有固定探测概率）、网络流分类（按协议类型拆分流量）。但若筛选概率依赖历史或到达时刻——例如系统繁忙时丢弃概率增大——两类流的独立性就不成立了。

\subsection{非齐次 Poisson 过程}

如果事件发生率本身随时间有规律地变化——比如网站流量白天高夜间低——可以用确定性函数 \(\lambda(t)\ge0\) 替代常数 \(\lambda\)。定义累计强度

\[
\Lambda(t)=\int_0^t\lambda(u)\,du.
\]

非齐次 Poisson 过程保留独立增量，但放弃平稳增量：

\[
N(t)-N(s)
\sim\mathrm{Poisson}
\left(\int_s^t\lambda(u)\,du\right).
\]

这里的 \(\lambda(t)\) 是预先给定的确定性函数，不相交区间仍有独立增量。但如果强度本身是随机的——例如依赖历史事件——就超出了非齐次 Poisson 的范围，进入 Cox 过程（随机强度）或 Hawkes 自激过程（事件触发强度上升）。不能直接沿用独立增量结论。

\subsection{怎么判断 Poisson 模型是否合适}

拿到数据后别急着拟合。先做几个检查：

\begin{itemize}
\tightlist
\item
  同长度窗口的计数均值是否随时间稳定？如果不稳定，考虑非齐次；
\item
  计数方差是否接近均值？方差远超均值提示过度离散，可能需要负二项或混合 Poisson；
\item
  不相交窗口是否近似独立？如果存在自相关，可能需要时间序列模型或自激过程；
\item
  间隔是否近似 Exponential？间隔的直方图形状是快速诊断的好工具；
\item
  是否出现过多零窗口（零膨胀）或成批到达（聚集效应）？
\end{itemize}

模型失效的方向本身会提示替代结构：时间变速率、混合 Poisson、更新过程、Hawkes 自激过程或有容量限制的过程。诊断的价值在于知道错在哪里才能找到该往哪里走。
